\documentclass[10pt,a4paper]{amsart}
\usepackage[margin=19mm]{geometry}
\usepackage{amsmath,amssymb,mathtools}
\usepackage[scr=rsfs,cal=euler]{mathalfa}
\usepackage{xfrac}
\usepackage[unicode,hidelinks,bookmarks=false]{hyperref}
\newcommand{\ie}{i.e.\ }
\newcommand{\cf}{cf.\ }
\newcommand{\resp}{respectively\ }
\newcommand{\Subsection}{\subsection}
\providecommand{\nobreakdash}{\nobreak\hbox{-}}
\setlength{\emergencystretch}{2em}
\multlinegap0pt
\usepackage{siunitx}
\usepackage{siunitx}
\usepackage{siunitx}
\usepackage{siunitx}
\usepackage{amssymb}
\usepackage{mathtools}
\usepackage{enumerate}
\usepackage{xcolor}
\usepackage{siunitx}
\newtheorem{remark}{Remark}
\newtheorem{lemma}{Lemma}
\newtheorem{theorem}{Theorem}
\newtheorem{corollary}{Corollary}
\usepackage{cleveref}
\crefname{remark}{Remark}{Remarks}
\crefname{lemma}{Lemma}{Lemmas}
\crefname{theorem}{Theorem}{Theorems}
\crefname{corollary}{Corollary}{Corollarys}
\newcommand{\height}{H}
\title{Completely Additive Height Functions: Profile Laws, Matula Bounds, and Inverse Growth}
\author{Hartosh Singh Bal}
\date{Source: arXiv:2308.00455v3 (CC BY 4.0)}
\begin{document}
\maketitle

\noindent\emph{Source and licence.} Completely Additive Height Functions: Profile Laws, Matula Bounds, and Inverse Growth. Version arXiv:2308.00455v3 (CC BY 4.0), distributed by arXiv under the Creative Commons Attribution 4.0 International licence, http://creativecommons.org/licenses/by/4.0/, as recorded in the arXiv licence metadata for this exact version and shown on the abstract page https://arxiv.org/abs/2308.00455v3. Full licence notice: \texttt{licenses/arxiv-2308.00455-CC-BY-4.0.txt}.\par\smallskip

\noindent\textit{Editorial scope.} Counted unit: the article's corollary cor:seq-avgorder-scale (source0.tex lines 2234--2253). The article's complete original proof of it is delivered with it (source0.tex lines 2255--2314, one \texttt{proof} environment of 60 lines). No mathematics is written, repaired or re-derived: the statement and the proof are the article's own lines, in the article's own notation, and no retained line is substituted or re-flowed.  Retained uncounted as the source-local context the proof needs: lines 1594--1604; lines 1646--1658; lines 2154--2161; lines 2193--2205.  Nothing was omitted from any retained span, so the package carries no omission record.  No retained line contains a citation, so the wrapper carries no bibliography and no bibliography entry.  No retained line contains a figure environment, an image-inclusion command or drawing code, so no image is delivered and the package declares no asset dependency.  Editorial formatting only, in addition: an AMS \texttt{amsart} preamble that restates the article's own declarations and macros used by the retained lines, namely the environments \texttt{remark}, \texttt{lemma}, \texttt{theorem}, \texttt{corollary} on separate counters, together with the standard packages those lines need.  The retained environments print in this document as Remark 1, Lemma 1, Theorem 1, Corollary 1 in order of appearance, whereas the article prints the same items with the numbering of its own full document.  The complete native source of the article as distributed in the arXiv e-print for this exact version is bundled unchanged as \texttt{sources/145066/source0.tex}.
\begin{equation}\label{eq:Fx-main-term}
  F(x)
  =
  x\sum_{p\le x}\frac{H(p)}{p-1}
  +
  O\!\left(
    \sum_{p\le x} H(p)\frac{\log x}{\log p}
  \right)
  +
  O\!\left(\sum_{p\le x}H(p)\right).
\end{equation}

\begin{remark}\label{rem:pminus1}
The replacement of $(p-1)^{-1}$ by $p^{-1}$ must be handled with some care.
In general
\[
  \frac1{p-1}=\frac1p+O\!\left(\frac1{p^2}\right),
\]
so the resulting change in the main sum is bounded by
\[
  x\sum_{p\le x}\frac{H(p)}{p^2}.
\]
This is lower order in the examples where we make the replacement, but it is not
uniformly $O(x)$ without additional hypotheses on $H(p)$.
\end{remark}

\begin{lemma}[A logarithmic-integral sum]\label{lem:logint-sum}
Fix $\beta>0$. Then, as $m\to\infty$,
\[
\sum_{2\le i\le m}\frac{i^{\beta-1}}{\log i}
\;=\;
\frac{m^\beta}{\beta\log m}\,(1+o(1)).
\]
\end{lemma}

\begin{theorem}[Prime-harmonic growth for the canonical sequential model]\label{thm:seq-prime-sum}
Assume $\Pi_k\sim Ck^\alpha$ with $C>0$ and $\alpha>0$, and let $m=\pi(x)$.
Then
\[
S_{\mathrm{seq}}(x)=\sum_{p\le x}\frac{\height_{\mathrm{seq}}(p)}{p}
\;=\;
\frac{\alpha}{C^{1/\alpha}}\cdot \frac{m^{1/\alpha}}{\log m}\,(1+o(1)).
\]
In particular, using $\pi(x)\sim x/\log x$,
\[
S_{\mathrm{seq}}(x)\asymp \frac{x^{1/\alpha}}{(\log x)^{1+1/\alpha}}.
\]
\end{theorem}

\begin{corollary}[Average-order scale from Equation~\eqref{eq:Fx-main-term}]
\label{cor:seq-avgorder-scale}
Assume $\Pi_k\sim Ck^\alpha$ as in \Cref{thm:seq-prime-sum}, and let
$F_{\mathrm{seq}}(x)=\sum_{n\le x}\height_{\mathrm{seq}}(n)$.  Then the leading
prime-sum term has scale
\[
  x\sum_{p\le x}\frac{\height_{\mathrm{seq}}(p)}{p}
  \asymp
  \frac{x^{1+1/\alpha}}{(\log x)^{1+1/\alpha}}.
\]
Moreover, the error term in Equation~\eqref{eq:Fx-main-term} has the same order of
magnitude.  Consequently
\[
  F_{\mathrm{seq}}(x)
  \asymp
  \frac{x^{1+1/\alpha}}{(\log x)^{1+1/\alpha}}.
\]
The reduction determines the order of magnitude in the polynomial-profile
regime, but not an asymptotic constant for $F_{\mathrm{seq}}(x)$.
\end{corollary}

\begin{proof}
Let $m=\pi(x)$.  By \Cref{thm:seq-prime-sum},
\[
  xS_{\mathrm{seq}}(x)
  \asymp
  x\frac{m^{1/\alpha}}{\log m}
  \asymp
  \frac{x^{1+1/\alpha}}{(\log x)^{1+1/\alpha}},
\]
using $m\sim x/\log x$.

For the error term in Equation~\eqref{eq:Fx-main-term}, use
$\height_{\mathrm{seq}}(p_i)\asymp i^{1/\alpha}$ and
$p_i\sim i\log i$.  Then
\[
\sum_{p\le x}\height_{\mathrm{seq}}(p)\frac{\log x}{\log p}
\asymp
\log x\sum_{i\le m}\frac{i^{1/\alpha}}{\log i}.
\]
Applying \Cref{lem:logint-sum} with $\beta=1+1/\alpha$ gives
\[
\sum_{i\le m}\frac{i^{1/\alpha}}{\log i}
\asymp
\frac{m^{1+1/\alpha}}{\log m}.
\]
Since $\log x\asymp\log m$, this error term is
\[
  \asymp m^{1+1/\alpha}
  \asymp
  \frac{x^{1+1/\alpha}}{(\log x)^{1+1/\alpha}}.
\]
For the lower bound, use the exact identity
\[
  F_{\mathrm{seq}}(x)
  =
  \sum_{p^\nu\le x}
  \height_{\mathrm{seq}}(p)
  \left\lfloor \frac{x}{p^\nu}\right\rfloor.
\]
Keeping only the terms with $\nu=1$ and $p\le x/2$ gives
\[
  F_{\mathrm{seq}}(x)
  \ge
  \sum_{p\le x/2}
  \height_{\mathrm{seq}}(p)
  \left\lfloor \frac{x}{p}\right\rfloor
  \ge
  \frac{x}{2}
  \sum_{p\le x/2}\frac{\height_{\mathrm{seq}}(p)}p,
\]
which has the same order of magnitude by \Cref{thm:seq-prime-sum}.

Finally, replacing $(p-1)^{-1}$ by $p^{-1}$ changes the leading prime sum by
\[
  x\sum_{p\le x}\frac{\height_{\mathrm{seq}}(p)}{p^2},
\]
which is lower order than $xS_{\mathrm{seq}}(x)$ by the same calculation as in
Remark~\ref{rem:pminus1}.  This proves the asserted order of magnitude and
explains why no asymptotic constant follows from the reduction alone.
\end{proof}

\end{document}
